\makeatletter
\def\input@path{{./}{../}{../../}{../../../}{../../../../}{preamble/}{../preamble/}{../../preamble/}{../../../preamble/}}
\makeatother

%\input{preamble_G10_T1_practice_activities_c1_c3.tex}
\input{preamble.tex}
\input{graphs/highlights.tex}
% \GTenHighlightsfalse % Uncomment to disable instructional highlights.
\input{graphs/specific_highlights.tex}


\title{10th Grade Midterm Analysis of Decisions: C4 Solutions}
\author{}
\date{}

\pagestyle{empty}

\begin{document}
	
	\maketitle
	\vspace*{-2.5cm}
	
	% ============================================================
	% CRITERION
	% ============================================================
	
	\section*{Evaluated Criterion}
	
	C4: Applies the Maximin and Maximax criteria to make optimistic or
	conservative decisions without probabilities, selecting the alternative
	with the highest possible payoff or the best worst-case outcome,
	interpreting risks or limitations.
	
	% ============================================================
	% INTRODUCTION
	% ============================================================
	
	\section{Maximax and Maximin Criteria}
\phantomsection\label{tone:legend}
\PracticeHighlightLegend{G10-T1-C4}

	
	When \GTenToneHL{probability}{probabilities} are not available, a decision-maker can compare
	\GTenToneHL{alternative}{alternatives} using the \GTenToneHL{payoff}{payoffs} associated with each possible state of
	nature.
	
	Two important decision criteria are:
	
	\begin{itemize}[leftmargin=*]
		\item \textbf{\GTenToneHL{best}{Maximax}}: an optimistic criterion that focuses on the
		best possible \GTenToneHL{payoff}{payoff} for each \GTenToneHL{alternative}{alternative}.
		
		\item \textbf{\GTenToneHL{worst}{Maximin}}: a conservative criterion that focuses on the
		worst possible \GTenToneHL{payoff}{payoff} for each \GTenToneHL{alternative}{alternative}.
	\end{itemize}
	
	For a \GTenToneHL{payoff}{payoff} table with \GTenToneHL{alternative}{alternatives} in rows and \GTenToneHL{state}{states} \GTenToneHL{state}{of} \GTenToneHL{state}{nature} in
	columns, let
	
	\[
	\GTenToneHL{payoff}{P_{ay}}(\GTenToneHL{alternative}{A_i},\GTenToneHL{state}{S_j})
	\]
	
	represent the \GTenToneHL{payoff}{payoff} obtained when \GTenToneHL{alternative}{alternative} $\GTenToneHL{alternative}{A_i}$ is selected and
	\GTenToneHL{state}{state} \GTenToneHL{state}{of} \GTenToneHL{state}{nature} $\GTenToneHL{state}{S_j}$ occurs.
	
	For each \GTenToneHL{alternative}{alternative}, calculate
	
	\[
	\GTenToneHL{best}{M_i}
	=
	\max_j \GTenToneHL{payoff}{P_{ay}}(\GTenToneHL{alternative}{A_i},\GTenToneHL{state}{S_j})
	=
	\max
	\left\{
	\GTenToneHL{payoff}{P_{ay}}(\GTenToneHL{alternative}{A_i},\GTenToneHL{state}{S_1}),
	\GTenToneHL{payoff}{P_{ay}}(\GTenToneHL{alternative}{A_i},\GTenToneHL{state}{S_2}),
	\ldots,
	\GTenToneHL{payoff}{P_{ay}}(\GTenToneHL{alternative}{A_i},\GTenToneHL{state}{S_n})
	\right\}
	\quad \longrightarrow \quad 
	A^*_{\mathrm{Maximax}}
	=
	\arg\max_i \GTenToneHL{best}{M_i}
	\]
	
	for the \GTenToneHL{best}{Maximax} criterion, and
	
	\[
	\GTenToneHL{worst}{m_i}
	=
	\min_j \GTenToneHL{payoff}{P_{ay}}(\GTenToneHL{alternative}{A_i},\GTenToneHL{state}{S_j})
	=
	\min
	\left\{
	\GTenToneHL{payoff}{P_{ay}}(\GTenToneHL{alternative}{A_i},\GTenToneHL{state}{S_1}),
	\GTenToneHL{payoff}{P_{ay}}(\GTenToneHL{alternative}{A_i},\GTenToneHL{state}{S_2}),
	\ldots,
	\GTenToneHL{payoff}{P_{ay}}(\GTenToneHL{alternative}{A_i},\GTenToneHL{state}{S_n})
	\right\}
	\quad \longrightarrow \quad
	A^*_{\mathrm{Maximin}}
	=
	\arg\max_i \GTenToneHL{worst}{m_i}
	\]
	
	for the \GTenToneHL{worst}{Maximin} criterion.
	
	% ============================================================
	% GENERAL 2x2 PROBLEM
	% ============================================================
	
	\section{General 2x2 Problem}
	
	Consider the following \GTenToneHL{payoff}{payoff} table:
	
	\begin{center}
		\renewcommand{\arraystretch}{1.3}
		\begin{tabular}{lcc}
			\toprule
			Alternative & $\GTenToneHL{state}{S_1}$ & $\GTenToneHL{state}{S_2}$ \\
			\midrule
			$\GTenToneHL{alternative}{A_1}$
			& $\GTenToneHL{payoff}{P_{ay}}(\GTenToneHL{alternative}{A_1},\GTenToneHL{state}{S_1})$
			& $\GTenToneHL{payoff}{P_{ay}}(\GTenToneHL{alternative}{A_1},\GTenToneHL{state}{S_2})$ \\
			
			$\GTenToneHL{alternative}{A_2}$
			& $\GTenToneHL{payoff}{P_{ay}}(\GTenToneHL{alternative}{A_2},\GTenToneHL{state}{S_1})$
			& $\GTenToneHL{payoff}{P_{ay}}(\GTenToneHL{alternative}{A_2},\GTenToneHL{state}{S_2})$ \\
			\bottomrule
		\end{tabular}
	\end{center}
	
	Determine the decision using both the \GTenToneHL{best}{Maximax} and \GTenToneHL{worst}{Maximin} criteria.
	
	% ============================================================
	% SOLUTION 2x2
	% ============================================================
	
	\subsection{Solution}
	
	\textbf{Step 1: Identify the \GTenToneHL{alternative}{alternatives} and \GTenToneHL{state}{states} \GTenToneHL{state}{of} \GTenToneHL{state}{nature}}
	
	Alternatives: $\GTenToneHL{alternative}{A_1}, \GTenToneHL{alternative}{A_2}$ \qquad \qquad States of nature: $\GTenToneHL{state}{S_1}, \GTenToneHL{state}{S_2}$.
	
	\noindent
	\setlength{\columnseprule}{0.4pt}
	
	\begin{multicols}{2}
		\vspace{0pt}
		
		
		
		\textbf{Step 2: Calculate the \GTenToneHL{best}{Maximax} values}
		
		For each \GTenToneHL{alternative}{alternative}, identify its highest possible \GTenToneHL{payoff}{payoff}:
		
		\[
		\GTenToneHL{best}{M_1}
		=
		\max
		\left\{
		\GTenToneHL{payoff}{P_{ay}}(\GTenToneHL{alternative}{A_1},\GTenToneHL{state}{S_1}),
		\GTenToneHL{payoff}{P_{ay}}(\GTenToneHL{alternative}{A_1},\GTenToneHL{state}{S_2})
		\right\},
		\]
		
		\[
		\GTenToneHL{best}{M_2}
		=
		\max
		\left\{
		\GTenToneHL{payoff}{P_{ay}}(\GTenToneHL{alternative}{A_2},\GTenToneHL{state}{S_1}),
		\GTenToneHL{payoff}{P_{ay}}(\GTenToneHL{alternative}{A_2},\GTenToneHL{state}{S_2})
		\right\}.
		\]
		
		Compare the two maximum values:
		
		\[
		A^*_{\mathrm{Maximax}}
		=
		\arg\max_i \GTenToneHL{best}{M_i}.
		\]
		
		Thus, choose the \GTenToneHL{alternative}{alternative} with the highest possible \GTenToneHL{payoff}{payoff}.
		
		\columnbreak
		
		\textbf{Step 3: Calculate the \GTenToneHL{worst}{Maximin} values}
		
		For each \GTenToneHL{alternative}{alternative}, identify its lowest possible \GTenToneHL{payoff}{payoff}:
		
		\[
		\GTenToneHL{worst}{m_1}
		=
		\min
		\left\{
		\GTenToneHL{payoff}{P_{ay}}(\GTenToneHL{alternative}{A_1},\GTenToneHL{state}{S_1}),
		\GTenToneHL{payoff}{P_{ay}}(\GTenToneHL{alternative}{A_1},\GTenToneHL{state}{S_2})
		\right\},
		\]
		
		\[
		\GTenToneHL{worst}{m_2}
		=
		\min
		\left\{
		\GTenToneHL{payoff}{P_{ay}}(\GTenToneHL{alternative}{A_2},\GTenToneHL{state}{S_1}),
		\GTenToneHL{payoff}{P_{ay}}(\GTenToneHL{alternative}{A_2},\GTenToneHL{state}{S_2})
		\right\}.
		\]
		
		Compare the two minimum values:
		
		\[
		A^*_{\mathrm{Maximin}}
		=
		\arg\max_i \GTenToneHL{worst}{m_i}.
		\]
		
		Thus, choose the \GTenToneHL{alternative}{alternative} with the highest worst-case \GTenToneHL{payoff}{payoff}.
		
	\end{multicols}
	
	\setlength{\columnseprule}{0pt}
	
	
	\textbf{Extended \GTenToneHL{payoff}{payoff} table}
	
	\begin{center}
		\small
		\renewcommand{\arraystretch}{1.5}
		\setlength{\tabcolsep}{3pt}
		
		\begin{tabular}{lcccc}
			\toprule
			Alternative & $\GTenToneHL{state}{S_1}$ & $\GTenToneHL{state}{S_2}$ & Max & Min \\
			\midrule
			
			$\GTenToneHL{alternative}{A_1}$
			& $\GTenToneHL{payoff}{P_{ay}}(\GTenToneHL{alternative}{A_1},\GTenToneHL{state}{S_1})$
			& $\GTenToneHL{payoff}{P_{ay}}(\GTenToneHL{alternative}{A_1},\GTenToneHL{state}{S_2})$
			& $\max\{\GTenToneHL{payoff}{P_{ay}}(\GTenToneHL{alternative}{A_1},\GTenToneHL{state}{S_1}),\GTenToneHL{payoff}{P_{ay}}(\GTenToneHL{alternative}{A_1},\GTenToneHL{state}{S_2})\}$
			& $\min\{\GTenToneHL{payoff}{P_{ay}}(\GTenToneHL{alternative}{A_1},\GTenToneHL{state}{S_1}),\GTenToneHL{payoff}{P_{ay}}(\GTenToneHL{alternative}{A_1},\GTenToneHL{state}{S_2})\}$
			\\
			
			$\GTenToneHL{alternative}{A_2}$
			& $\GTenToneHL{payoff}{P_{ay}}(\GTenToneHL{alternative}{A_2},\GTenToneHL{state}{S_1})$
			& $\GTenToneHL{payoff}{P_{ay}}(\GTenToneHL{alternative}{A_2},\GTenToneHL{state}{S_2})$
			& $\max\{\GTenToneHL{payoff}{P_{ay}}(\GTenToneHL{alternative}{A_2},\GTenToneHL{state}{S_1}),\GTenToneHL{payoff}{P_{ay}}(\GTenToneHL{alternative}{A_2},\GTenToneHL{state}{S_2})\}$
			& $\min\{\GTenToneHL{payoff}{P_{ay}}(\GTenToneHL{alternative}{A_2},\GTenToneHL{state}{S_1}),\GTenToneHL{payoff}{P_{ay}}(\GTenToneHL{alternative}{A_2},\GTenToneHL{state}{S_2})\}$
			\\
			
			\bottomrule
		\end{tabular}
	\end{center}
	
	% ============================================================
	% GENERAL 4x3 PROBLEM
	% ============================================================
	
	\section{General 4x3 Problem}
	
		Consider the following \GTenToneHL{payoff}{payoff} table:
		
			\begin{center}
			\renewcommand{\arraystretch}{1.3}
			\begin{tabular}{lccc}
			\toprule
			Alternative & $\GTenToneHL{state}{S_1}$ & $\GTenToneHL{state}{S_2}$ & $\GTenToneHL{state}{S_3}$ \\
			\midrule
			
			$\GTenToneHL{alternative}{A_1}$
			& $\GTenToneHL{payoff}{P_{ay}}(\GTenToneHL{alternative}{A_1},\GTenToneHL{state}{S_1})$
			& $\GTenToneHL{payoff}{P_{ay}}(\GTenToneHL{alternative}{A_1},\GTenToneHL{state}{S_2})$
			& $\GTenToneHL{payoff}{P_{ay}}(\GTenToneHL{alternative}{A_1},\GTenToneHL{state}{S_3})$
			\\
			
			$\GTenToneHL{alternative}{A_2}$
			& $\GTenToneHL{payoff}{P_{ay}}(\GTenToneHL{alternative}{A_2},\GTenToneHL{state}{S_1})$
			& $\GTenToneHL{payoff}{P_{ay}}(\GTenToneHL{alternative}{A_2},\GTenToneHL{state}{S_2})$
			& $\GTenToneHL{payoff}{P_{ay}}(\GTenToneHL{alternative}{A_2},\GTenToneHL{state}{S_3})$
			\\
			
			$\GTenToneHL{alternative}{A_3}$
			& $\GTenToneHL{payoff}{P_{ay}}(\GTenToneHL{alternative}{A_3},\GTenToneHL{state}{S_1})$
			& $\GTenToneHL{payoff}{P_{ay}}(\GTenToneHL{alternative}{A_3},\GTenToneHL{state}{S_2})$
			& $\GTenToneHL{payoff}{P_{ay}}(\GTenToneHL{alternative}{A_3},\GTenToneHL{state}{S_3})$
			\\
			
			$\GTenToneHL{alternative}{A_4}$
			& $\GTenToneHL{payoff}{P_{ay}}(\GTenToneHL{alternative}{A_4},\GTenToneHL{state}{S_1})$
			& $\GTenToneHL{payoff}{P_{ay}}(\GTenToneHL{alternative}{A_4},\GTenToneHL{state}{S_2})$
			& $\GTenToneHL{payoff}{P_{ay}}(\GTenToneHL{alternative}{A_4},\GTenToneHL{state}{S_3})$
		\\
	
	\bottomrule
	\end{tabular}
	\end{center}
	
	Determine the decision using both the \GTenToneHL{best}{Maximax} and \GTenToneHL{worst}{Maximin} criteria.
	
	% ============================================================
	% SOLUTION 4x3
	% ============================================================
	
	\subsection{Solution}
	
	\textbf{Step 1: Identify the \GTenToneHL{alternative}{alternatives} and \GTenToneHL{state}{states} \GTenToneHL{state}{of} \GTenToneHL{state}{nature}}
	
	Alternatives: $\GTenToneHL{alternative}{A_1}, \GTenToneHL{alternative}{A_2}, \GTenToneHL{alternative}{A_3}, \GTenToneHL{alternative}{A_4}$
	\qquad
	States of nature: $\GTenToneHL{state}{S_1}, \GTenToneHL{state}{S_2}, \GTenToneHL{state}{S_3}$.
	
	\noindent
		\setlength{\columnseprule}{0.4pt}
		
		\begin{multicols}{2}
		\vspace{0pt}
		
		\textbf{Step 2: Calculate the \GTenToneHL{best}{Maximax} values}
		
		For each \GTenToneHL{alternative}{alternative}, identify its highest possible \GTenToneHL{payoff}{payoff}:
		
		\[
		\GTenToneHL{best}{M_1}
		=
		\max
		\left\{
		\GTenToneHL{payoff}{P_{ay}}(\GTenToneHL{alternative}{A_1},\GTenToneHL{state}{S_1}),
		\GTenToneHL{payoff}{P_{ay}}(\GTenToneHL{alternative}{A_1},\GTenToneHL{state}{S_2}),
		\GTenToneHL{payoff}{P_{ay}}(\GTenToneHL{alternative}{A_1},\GTenToneHL{state}{S_3})
		\right\},
		\]
		
		\[
		\GTenToneHL{best}{M_2}
		=
		\max
		\left\{
		\GTenToneHL{payoff}{P_{ay}}(\GTenToneHL{alternative}{A_2},\GTenToneHL{state}{S_1}),
		\GTenToneHL{payoff}{P_{ay}}(\GTenToneHL{alternative}{A_2},\GTenToneHL{state}{S_2}),
		\GTenToneHL{payoff}{P_{ay}}(\GTenToneHL{alternative}{A_2},\GTenToneHL{state}{S_3})
		\right\},
		\]
		
		\[
		\GTenToneHL{best}{M_3}
		=
		\max
		\left\{
		\GTenToneHL{payoff}{P_{ay}}(\GTenToneHL{alternative}{A_3},\GTenToneHL{state}{S_1}),
		\GTenToneHL{payoff}{P_{ay}}(\GTenToneHL{alternative}{A_3},\GTenToneHL{state}{S_2}),
		\GTenToneHL{payoff}{P_{ay}}(\GTenToneHL{alternative}{A_3},\GTenToneHL{state}{S_3})
		\right\},
		\]
		
		\[
		\GTenToneHL{best}{M_4}
		=
		\max
		\left\{
		\GTenToneHL{payoff}{P_{ay}}(\GTenToneHL{alternative}{A_4},\GTenToneHL{state}{S_1}),
		\GTenToneHL{payoff}{P_{ay}}(\GTenToneHL{alternative}{A_4},\GTenToneHL{state}{S_2}),
		\GTenToneHL{payoff}{P_{ay}}(\GTenToneHL{alternative}{A_4},\GTenToneHL{state}{S_3})
		\right\}.
		\]
		
		Compare the four maximum values:
		
		\[
		A^*_{\mathrm{Maximax}}
		=
		\arg\max_i \GTenToneHL{best}{M_i}.
		\]
		
		Thus, choose the \GTenToneHL{alternative}{alternative} with the highest possible \GTenToneHL{payoff}{payoff}.
		
		\columnbreak
		
		\textbf{Step 3: Calculate the \GTenToneHL{worst}{Maximin} values}
		
		For each \GTenToneHL{alternative}{alternative}, identify its lowest possible \GTenToneHL{payoff}{payoff}:
		
		\[
		\GTenToneHL{worst}{m_1}
		=
		\min
		\left\{
		\GTenToneHL{payoff}{P_{ay}}(\GTenToneHL{alternative}{A_1},\GTenToneHL{state}{S_1}),
		\GTenToneHL{payoff}{P_{ay}}(\GTenToneHL{alternative}{A_1},\GTenToneHL{state}{S_2}),
		\GTenToneHL{payoff}{P_{ay}}(\GTenToneHL{alternative}{A_1},\GTenToneHL{state}{S_3})
		\right\},
		\]
		
		\[
		\GTenToneHL{worst}{m_2}
		=
		\min
		\left\{
		\GTenToneHL{payoff}{P_{ay}}(\GTenToneHL{alternative}{A_2},\GTenToneHL{state}{S_1}),
		\GTenToneHL{payoff}{P_{ay}}(\GTenToneHL{alternative}{A_2},\GTenToneHL{state}{S_2}),
		\GTenToneHL{payoff}{P_{ay}}(\GTenToneHL{alternative}{A_2},\GTenToneHL{state}{S_3})
		\right\},
		\]
		
		\[
		\GTenToneHL{worst}{m_3}
		=
		\min
		\left\{
		\GTenToneHL{payoff}{P_{ay}}(\GTenToneHL{alternative}{A_3},\GTenToneHL{state}{S_1}),
		\GTenToneHL{payoff}{P_{ay}}(\GTenToneHL{alternative}{A_3},\GTenToneHL{state}{S_2}),
		\GTenToneHL{payoff}{P_{ay}}(\GTenToneHL{alternative}{A_3},\GTenToneHL{state}{S_3})
		\right\},
		\]
		
		\[
		\GTenToneHL{worst}{m_4}
		=
		\min
		\left\{
		\GTenToneHL{payoff}{P_{ay}}(\GTenToneHL{alternative}{A_4},\GTenToneHL{state}{S_1}),
		\GTenToneHL{payoff}{P_{ay}}(\GTenToneHL{alternative}{A_4},\GTenToneHL{state}{S_2}),
		\GTenToneHL{payoff}{P_{ay}}(\GTenToneHL{alternative}{A_4},\GTenToneHL{state}{S_3})
		\right\}.
		\]
		
		Compare the four minimum values:
		
		\[
		A^*_{\mathrm{Maximin}}
		=
		\arg\max_i \GTenToneHL{worst}{m_i}.
		\]
	
	Thus, choose the \GTenToneHL{alternative}{alternative} with the highest worst-case \GTenToneHL{payoff}{payoff}.
	
	\end{multicols}
	
	\setlength{\columnseprule}{0pt}
	
		\textbf{Extended \GTenToneHL{payoff}{payoff} table}
		
		\begin{center}
		\scriptsize
		\renewcommand{\arraystretch}{1.5}
			\setlength{\tabcolsep}{2pt}
			
			\begin{tabular}{lccccc}
			\toprule
			Alternative & $\GTenToneHL{state}{S_1}$ & $\GTenToneHL{state}{S_2}$ & $\GTenToneHL{state}{S_3}$ & Max & Min \\
			\midrule
			
			$\GTenToneHL{alternative}{A_1}$
			& $\GTenToneHL{payoff}{P_{ay}}(\GTenToneHL{alternative}{A_1},\GTenToneHL{state}{S_1})$
			& $\GTenToneHL{payoff}{P_{ay}}(\GTenToneHL{alternative}{A_1},\GTenToneHL{state}{S_2})$
			& $\GTenToneHL{payoff}{P_{ay}}(\GTenToneHL{alternative}{A_1},\GTenToneHL{state}{S_3})$
			& $\max\{\GTenToneHL{payoff}{P_{ay}}(\GTenToneHL{alternative}{A_1},\GTenToneHL{state}{S_1}),\GTenToneHL{payoff}{P_{ay}}(\GTenToneHL{alternative}{A_1},\GTenToneHL{state}{S_2}),\GTenToneHL{payoff}{P_{ay}}(\GTenToneHL{alternative}{A_1},\GTenToneHL{state}{S_3})\}$
			& $\min\{\GTenToneHL{payoff}{P_{ay}}(\GTenToneHL{alternative}{A_1},\GTenToneHL{state}{S_1}),\GTenToneHL{payoff}{P_{ay}}(\GTenToneHL{alternative}{A_1},\GTenToneHL{state}{S_2}),\GTenToneHL{payoff}{P_{ay}}(\GTenToneHL{alternative}{A_1},\GTenToneHL{state}{S_3})\}$
			\\
			
			$\GTenToneHL{alternative}{A_2}$
			& $\GTenToneHL{payoff}{P_{ay}}(\GTenToneHL{alternative}{A_2},\GTenToneHL{state}{S_1})$
			& $\GTenToneHL{payoff}{P_{ay}}(\GTenToneHL{alternative}{A_2},\GTenToneHL{state}{S_2})$
			& $\GTenToneHL{payoff}{P_{ay}}(\GTenToneHL{alternative}{A_2},\GTenToneHL{state}{S_3})$
			& $\max\{\GTenToneHL{payoff}{P_{ay}}(\GTenToneHL{alternative}{A_2},\GTenToneHL{state}{S_1}),\GTenToneHL{payoff}{P_{ay}}(\GTenToneHL{alternative}{A_2},\GTenToneHL{state}{S_2}),\GTenToneHL{payoff}{P_{ay}}(\GTenToneHL{alternative}{A_2},\GTenToneHL{state}{S_3})\}$
			& $\min\{\GTenToneHL{payoff}{P_{ay}}(\GTenToneHL{alternative}{A_2},\GTenToneHL{state}{S_1}),\GTenToneHL{payoff}{P_{ay}}(\GTenToneHL{alternative}{A_2},\GTenToneHL{state}{S_2}),\GTenToneHL{payoff}{P_{ay}}(\GTenToneHL{alternative}{A_2},\GTenToneHL{state}{S_3})\}$
			\\
			
			$\GTenToneHL{alternative}{A_3}$
			& $\GTenToneHL{payoff}{P_{ay}}(\GTenToneHL{alternative}{A_3},\GTenToneHL{state}{S_1})$
			& $\GTenToneHL{payoff}{P_{ay}}(\GTenToneHL{alternative}{A_3},\GTenToneHL{state}{S_2})$
			& $\GTenToneHL{payoff}{P_{ay}}(\GTenToneHL{alternative}{A_3},\GTenToneHL{state}{S_3})$
			& $\max\{\GTenToneHL{payoff}{P_{ay}}(\GTenToneHL{alternative}{A_3},\GTenToneHL{state}{S_1}),\GTenToneHL{payoff}{P_{ay}}(\GTenToneHL{alternative}{A_3},\GTenToneHL{state}{S_2}),\GTenToneHL{payoff}{P_{ay}}(\GTenToneHL{alternative}{A_3},\GTenToneHL{state}{S_3})\}$
			& $\min\{\GTenToneHL{payoff}{P_{ay}}(\GTenToneHL{alternative}{A_3},\GTenToneHL{state}{S_1}),\GTenToneHL{payoff}{P_{ay}}(\GTenToneHL{alternative}{A_3},\GTenToneHL{state}{S_2}),\GTenToneHL{payoff}{P_{ay}}(\GTenToneHL{alternative}{A_3},\GTenToneHL{state}{S_3})\}$
			\\
			
			$\GTenToneHL{alternative}{A_4}$
			& $\GTenToneHL{payoff}{P_{ay}}(\GTenToneHL{alternative}{A_4},\GTenToneHL{state}{S_1})$
			& $\GTenToneHL{payoff}{P_{ay}}(\GTenToneHL{alternative}{A_4},\GTenToneHL{state}{S_2})$
			& $\GTenToneHL{payoff}{P_{ay}}(\GTenToneHL{alternative}{A_4},\GTenToneHL{state}{S_3})$
			& $\max\{\GTenToneHL{payoff}{P_{ay}}(\GTenToneHL{alternative}{A_4},\GTenToneHL{state}{S_1}),\GTenToneHL{payoff}{P_{ay}}(\GTenToneHL{alternative}{A_4},\GTenToneHL{state}{S_2}),\GTenToneHL{payoff}{P_{ay}}(\GTenToneHL{alternative}{A_4},\GTenToneHL{state}{S_3})\}$
			& $\min\{\GTenToneHL{payoff}{P_{ay}}(\GTenToneHL{alternative}{A_4},\GTenToneHL{state}{S_1}),\GTenToneHL{payoff}{P_{ay}}(\GTenToneHL{alternative}{A_4},\GTenToneHL{state}{S_2}),\GTenToneHL{payoff}{P_{ay}}(\GTenToneHL{alternative}{A_4},\GTenToneHL{state}{S_3})\}$
		\\
	
	\bottomrule
	\end{tabular}
	\end{center}
	
	
	% ============================================================ 
	% COMPARISON 
	% ============================================================
	
	
\phantomsection\label{tone:reference}
\subsection*{Notation reference}
\GTenToneReference{C4}
\section{Comparing the Two Criteria}
	
	The two criteria use the same \GTenToneHL{payoff}{payoff} table but emphasize different
	aspects of the decision.
	\begin{center}
		\renewcommand{\arraystretch}{1.5}
		\begin{tabular}{p{0.20\textwidth}p{0.32\textwidth}p{0.32\textwidth}}
			\toprule
			& \textbf{\GTenToneHL{best}{Maximax}} & \textbf{\GTenToneHL{worst}{Maximin}} \\
			\midrule
			
			Focus
			&
			Best possible \GTenToneHL{payoff}{payoff}
			&
			Worst possible \GTenToneHL{payoff}{payoff}
			\\
			
			Decision rule
			&
			Choose the \GTenToneHL{alternative}{alternative} with the highest possible \GTenToneHL{payoff}{payoff}
			&
			Choose the \GTenToneHL{alternative}{alternative} with the highest worst-case \GTenToneHL{payoff}{payoff}
			\\
			
			Attitude
			&
			Optimistic
			&
			Conservative
			\\
			
			Risk
			&
			Accepts greater uncertainty in pursuit of a high \GTenToneHL{payoff}{payoff}
			&
			Places greater emphasis on protection against poor outcomes
			\\
			
			Use when
			&
			The decision-maker focuses on favorable possibilities
			&
			The decision-maker wants protection against unfavorable possibilities
			\\
			
			\bottomrule
		\end{tabular}
	\end{center}
	
	A decision-maker using \textbf{\GTenToneHL{best}{Maximax}} asks:
	\[
	\text{``What is the best that could happen?''}
	\]
	A decision-maker using \textbf{\GTenToneHL{worst}{Maximin}} asks:
	\[
	\text{``What is the worst that could happen, and which alternative
		handles it best?''}
	\]
	
	\section{Contextualized 2x2 Problem - Coffee shop}
	
	A small coffee shop is considering whether to introduce a new dessert line or keep its current menu. The manager knows that the level of customer traffic during the promotional weekend is uncertain and has no reliable information about the \GTenToneHL{probability}{probability} of high or low traffic.
	
	What would a perfectly logical and optimistic decision-maker choose? What would a perfectly logical and pessimistic decision-maker choose?
	
	\begin{center}
		\renewcommand{\arraystretch}{1.3}
		\begin{tabular}{lcc}
			\toprule
			Alternative & High traffic ($\GTenToneHL{state}{S_1}$) & Low traffic ($\GTenToneHL{state}{S_2}$) \\
			\midrule
			Launch dessert line ($\GTenToneHL{alternative}{A_1}$) & \GTenToneHL{payoff}{42} & \GTenToneHL{payoff}{8} \\
			Keep current menu ($\GTenToneHL{alternative}{A_2}$) & \GTenToneHL{payoff}{28} & \GTenToneHL{payoff}{20} \\
			\bottomrule
		\end{tabular}
	\end{center}
	
	All \GTenToneHL{payoff}{payoffs} are expressed in thousands of dollars.
	
	Determine the decision using both the \GTenToneHL{best}{Maximax} and \GTenToneHL{worst}{Maximin} criteria.
	
	% ============================================================
	% SOLUTION 2x2
	% ============================================================
	
	\subsection{Solution}
	
	\textbf{Step 1: Identify the \GTenToneHL{alternative}{alternatives} and \GTenToneHL{state}{states} \GTenToneHL{state}{of} \GTenToneHL{state}{nature}}
	
	Alternatives: $\GTenToneHL{alternative}{A_1}$ = Launch dessert line, $\GTenToneHL{alternative}{A_2}$ = Keep current menu
	\qquad
	States of nature: $\GTenToneHL{state}{S_1}$ = High traffic, $\GTenToneHL{state}{S_2}$ = Low traffic.
	
	\noindent
	\setlength{\columnseprule}{0.4pt}
	
	\begin{multicols}{2}
		\vspace{0pt}
		
		\textbf{Step 2: Calculate the \GTenToneHL{best}{Maximax} values}
		
		For each \GTenToneHL{alternative}{alternative}, identify its highest possible \GTenToneHL{payoff}{payoff}:
		
		\[
		\GTenToneHL{best}{M_1}
		=
		\max\{42,8\}
		=
		42,
		\]
		
		\[
		\GTenToneHL{best}{M_2}
		=
		\max\{28,20\}
		=
		28.
		\]
		
		Compare the two maximum values:
		
		\[
		A^*_{\mathrm{Maximax}}
		=
		\arg\max_i \GTenToneHL{best}{M_i}
		=
		\GTenToneHL{alternative}{A_1}.
		\]
		
		Thus, the \GTenToneHL{best}{Maximax} criterion recommends launching the dessert line because it offers the highest possible \GTenToneHL{payoff}{payoff}.
		
		\columnbreak
		
		\textbf{Step 3: Calculate the \GTenToneHL{worst}{Maximin} values}
		
		For each \GTenToneHL{alternative}{alternative}, identify its lowest possible \GTenToneHL{payoff}{payoff}:
		
		\[
		\GTenToneHL{worst}{m_1}
		=
		\min\{42,8\}
		=
		8,
		\]
		
		\[
		\GTenToneHL{worst}{m_2}
		=
		\min\{28,20\}
		=
		20.
		\]
		
		Compare the two minimum values:
		
		\[
		A^*_{\mathrm{Maximin}}
		=
		\arg\max_i \GTenToneHL{worst}{m_i}
		=
		\GTenToneHL{alternative}{A_2}.
		\]
		
		Thus, the \GTenToneHL{worst}{Maximin} criterion recommends keeping the current menu because it provides the highest worst-case \GTenToneHL{payoff}{payoff}.
		
	\end{multicols}
	
	\setlength{\columnseprule}{0pt}
	
	\textbf{Extended \GTenToneHL{payoff}{payoff} table}
	
	\begin{center}
		\small
		\renewcommand{\arraystretch}{1.5}
		\setlength{\tabcolsep}{4pt}
		
		\begin{tabular}{lcccc}
			\toprule
			Alternative & $\GTenToneHL{state}{S_1}$ & $\GTenToneHL{state}{S_2}$ & Max & Min \\
			\midrule
			
			$\GTenToneHL{alternative}{A_1}$: Launch dessert line
			& $42$
			& $8$
			& $\max\{42,8\}=42$
			& $\min\{42,8\}=8$
			\\
			
			$\GTenToneHL{alternative}{A_2}$: Keep current menu
			& $28$
			& $20$
			& $\max\{28,20\}=28$
			& $\min\{28,20\}=20$
			\\
			
			\bottomrule
		\end{tabular}
		\end{center}
	
	\section{Contextualized 2x2 Problem - Clothing store}
	
	A local clothing store is deciding whether to launch a new advertising campaign before the holiday season or continue with its current marketing strategy. The manager knows that customer response to the campaign is uncertain and has no reliable information about the \GTenToneHL{probability}{probability} of strong or weak sales.
	
	What would a perfectly logical and optimistic decision-maker choose? What would a perfectly logical and pessimistic decision-maker choose?
	
	\begin{center}
		\renewcommand{\arraystretch}{1.3}
		\begin{tabular}{lcc}
			\toprule
			Alternative & Strong sales ($\GTenToneHL{state}{S_1}$) & Weak sales ($\GTenToneHL{state}{S_2}$) \\
			\midrule
			Launch new campaign ($\GTenToneHL{alternative}{A_1}$) & \GTenToneHL{payoff}{55} & \GTenToneHL{payoff}{5} \\
			Keep current strategy ($\GTenToneHL{alternative}{A_2}$) & \GTenToneHL{payoff}{35} & \GTenToneHL{payoff}{22} \\
			\bottomrule
		\end{tabular}
	\end{center}
	
	All \GTenToneHL{payoff}{payoffs} are expressed in thousands of dollars.
	
	Determine the decision using both the \GTenToneHL{best}{Maximax} and \GTenToneHL{worst}{Maximin} criteria.
	
	% ============================================================
	% SOLUTION 2x2
	% ============================================================
	
	\subsection{Solution}
	
	\textbf{Step 1: Identify the \GTenToneHL{alternative}{alternatives} and \GTenToneHL{state}{states} \GTenToneHL{state}{of} \GTenToneHL{state}{nature}}
	
	Alternatives: $\GTenToneHL{alternative}{A_1}$ = Launch new campaign, $\GTenToneHL{alternative}{A_2}$ = Keep current strategy
	\qquad
	States of nature: $\GTenToneHL{state}{S_1}$ = Strong sales, $\GTenToneHL{state}{S_2}$ = Weak sales.
	
	\noindent
	\setlength{\columnseprule}{0.4pt}
	
	\begin{multicols}{2}
		\vspace{0pt}
		
		\textbf{Step 2: Calculate the \GTenToneHL{best}{Maximax} values}
		
		For each \GTenToneHL{alternative}{alternative}, identify its highest possible \GTenToneHL{payoff}{payoff}:
		
		\[
		\GTenToneHL{best}{M_1}
		=
		\max\{55,5\}
		=
		55,
		\]
		
		\[
		\GTenToneHL{best}{M_2}
		=
		\max\{35,22\}
		=
		35.
		\]
		
		Compare the two maximum values:
		
		\[
		A^*_{\mathrm{Maximax}}
		=
		\arg\max_i \GTenToneHL{best}{M_i}
		=
		\GTenToneHL{alternative}{A_1}.
		\]
		
		Thus, the \GTenToneHL{best}{Maximax} criterion recommends launching the new campaign because it offers the highest possible \GTenToneHL{payoff}{payoff}.
		
		\columnbreak
		
		\textbf{Step 3: Calculate the \GTenToneHL{worst}{Maximin} values}
		
		For each \GTenToneHL{alternative}{alternative}, identify its lowest possible \GTenToneHL{payoff}{payoff}:
		
		\[
		\GTenToneHL{worst}{m_1}
		=
		\min\{55,5\}
		=
		5,
		\]
		
		\[
		\GTenToneHL{worst}{m_2}
		=
		\min\{35,22\}
		=
		22.
		\]
		
		Compare the two minimum values:
		
		\[
		A^*_{\mathrm{Maximin}}
		=
		\arg\max_i \GTenToneHL{worst}{m_i}
		=
		\GTenToneHL{alternative}{A_2}.
		\]
		
		Thus, the \GTenToneHL{worst}{Maximin} criterion recommends keeping the current strategy because it provides the highest worst-case \GTenToneHL{payoff}{payoff}.
		
	\end{multicols}
	
	\setlength{\columnseprule}{0pt}
	
	\textbf{Extended \GTenToneHL{payoff}{payoff} table}
	
	\begin{center}
		\small
		\renewcommand{\arraystretch}{1.5}
		\setlength{\tabcolsep}{4pt}
		
		\begin{tabular}{lcccc}
			\toprule
			Alternative & $\GTenToneHL{state}{S_1}$ & $\GTenToneHL{state}{S_2}$ & Max & Min \\
			\midrule
			
			$\GTenToneHL{alternative}{A_1}$: Launch new campaign
			& $55$
			& $5$
			& $\max\{55,5\}=55$
			& $\min\{55,5\}=5$
			\\
			
			$\GTenToneHL{alternative}{A_2}$: Keep current strategy
			& $35$
			& $22$
			& $\max\{35,22\}=35$
			& $\min\{35,22\}=22$
			\\
			
			\bottomrule
		\end{tabular}
	\end{center}
	
	
	\section{Contextualized 4x3 Problem - Technology company}
	
	A technology company is deciding which product strategy to adopt for the launch of a new device. The manager knows that market conditions are uncertain and has no reliable information about the \GTenToneHL{probability}{probability} of each possible level of demand. The company is considering four possible strategies: launching a premium version, launching a standard version, launching a budget version, or delaying the launch.
	
	The level of market demand could be high, medium, or low. The estimated \GTenToneHL{payoff}{profit} for each strategy under each possible market condition is shown in the table below.
	
	What would a perfectly logical and optimistic decision-maker choose? What would a perfectly logical and pessimistic decision-maker choose?
	
	\begin{center}
		\renewcommand{\arraystretch}{1.3}
		\begin{tabular}{lccc}
			\toprule
			Alternative & High demand ($\GTenToneHL{state}{S_1}$) & Medium demand ($\GTenToneHL{state}{S_2}$) & Low demand ($\GTenToneHL{state}{S_3}$) \\
			\midrule
			Premium version ($\GTenToneHL{alternative}{A_1}$) & \GTenToneHL{payoff}{80} & \GTenToneHL{payoff}{35} & \GTenToneHL{payoff}{-20} \\
			Standard version ($\GTenToneHL{alternative}{A_2}$) & \GTenToneHL{payoff}{60} & \GTenToneHL{payoff}{42} & \GTenToneHL{payoff}{18} \\
			Budget version ($\GTenToneHL{alternative}{A_3}$) & \GTenToneHL{payoff}{40} & \GTenToneHL{payoff}{32} & \GTenToneHL{payoff}{25} \\
			Delay launch ($\GTenToneHL{alternative}{A_4}$) & \GTenToneHL{payoff}{15} & \GTenToneHL{payoff}{15} & \GTenToneHL{payoff}{15} \\
			\bottomrule
		\end{tabular}
	\end{center}
	
	All \GTenToneHL{payoff}{payoffs} are expressed in thousands of dollars.
	
	Determine the decision using both the \GTenToneHL{best}{Maximax} and \GTenToneHL{worst}{Maximin} criteria.
	
	% ============================================================
	% SOLUTION 4x3
	% ============================================================
	
	\subsection{Solution}
	
	\textbf{Step 1: Identify the \GTenToneHL{alternative}{alternatives} and \GTenToneHL{state}{states} \GTenToneHL{state}{of} \GTenToneHL{state}{nature}}
	
	Alternatives: $\GTenToneHL{alternative}{A_1}$ = Premium version, $\GTenToneHL{alternative}{A_2}$ = Standard version,
	$\GTenToneHL{alternative}{A_3}$ = Budget version, $\GTenToneHL{alternative}{A_4}$ = Delay launch
	
	\qquad
	
	States of nature: $\GTenToneHL{state}{S_1}$ = High demand, $\GTenToneHL{state}{S_2}$ = Medium demand,
	$\GTenToneHL{state}{S_3}$ = Low demand.
	
	\noindent
	\setlength{\columnseprule}{0.4pt}
	
	\begin{multicols}{2}
		\vspace{0pt}
		
		\textbf{Step 2: Calculate the \GTenToneHL{best}{Maximax} values}
		
		For each \GTenToneHL{alternative}{alternative}, identify its highest possible \GTenToneHL{payoff}{payoff}:
		
		\[
		\GTenToneHL{best}{M_1}
		=
		\max\{80,35,-20\}
		=
		80,
		\]
		
		\[
		\GTenToneHL{best}{M_2}
		=
		\max\{60,42,18\}
		=
		60,
		\]
		
		\[
		\GTenToneHL{best}{M_3}
		=
		\max\{40,32,25\}
		=
		40,
		\]
		
		\[
		\GTenToneHL{best}{M_4}
		=
		\max\{15,15,15\}
		=
		15.
		\]
		
		Compare the four maximum values:
		
		\[
		A^*_{\mathrm{Maximax}}
		=
		\arg\max_i \GTenToneHL{best}{M_i}
		=
		\GTenToneHL{alternative}{A_1}.
		\]
		
		Thus, the \GTenToneHL{best}{Maximax} criterion recommends launching the premium version because it offers the highest possible \GTenToneHL{payoff}{payoff}.
		
		\columnbreak
		
		\textbf{Step 3: Calculate the \GTenToneHL{worst}{Maximin} values}
		
		For each \GTenToneHL{alternative}{alternative}, identify its lowest possible \GTenToneHL{payoff}{payoff}:
		
		\[
		\GTenToneHL{worst}{m_1}
		=
		\min\{80,35,-20\}
		=
		-20,
		\]
		
		\[
		\GTenToneHL{worst}{m_2}
		=
		\min\{60,42,18\}
		=
		18,
		\]
		
		\[
		\GTenToneHL{worst}{m_3}
		=
		\min\{40,32,25\}
		=
		25,
		\]
		
		\[
		\GTenToneHL{worst}{m_4}
		=
		\min\{15,15,15\}
		=
		15.
		\]
		
		Compare the four minimum values:
		
		\[
		A^*_{\mathrm{Maximin}}
		=
		\arg\max_i \GTenToneHL{worst}{m_i}
		=
		\GTenToneHL{alternative}{A_3}.
		\]
		
		Thus, the \GTenToneHL{worst}{Maximin} criterion recommends launching the budget version because it provides the highest worst-case \GTenToneHL{payoff}{payoff}.
		
	\end{multicols}
	
	\setlength{\columnseprule}{0pt}
	
	\textbf{Extended \GTenToneHL{payoff}{payoff} table}
	
	\begin{center}
		\scriptsize
		\renewcommand{\arraystretch}{1.5}
		\setlength{\tabcolsep}{3pt}
		
		\begin{tabular}{lccccc}
			\toprule
			Alternative & $\GTenToneHL{state}{S_1}$ & $\GTenToneHL{state}{S_2}$ & $\GTenToneHL{state}{S_3}$ & Max & Min \\
			\midrule
			
			$\GTenToneHL{alternative}{A_1}$: Premium version
			& $80$ & $35$ & $-20$
			& $\max\{80,35,-20\}=80$
			& $\min\{80,35,-20\}=-20$
			\\
			
			$\GTenToneHL{alternative}{A_2}$: Standard version
			& $60$ & $42$ & $18$
			& $\max\{60,42,18\}=60$
			& $\min\{60,42,18\}=18$
			\\
			
			$\GTenToneHL{alternative}{A_3}$: Budget version
			& $40$ & $32$ & $25$
			& $\max\{40,32,25\}=40$
			& $\min\{40,32,25\}=25$
			\\
			
			$\GTenToneHL{alternative}{A_4}$: Delay launch
			& $15$ & $15$ & $15$
			& $\max\{15,15,15\}=15$
			& $\min\{15,15,15\}=15$
			\\
			
			\bottomrule
		\end{tabular}
	\end{center}
	
	
	
\end{document}
